% ============================================================
\section{Related Work}
\label{sec:related}
% ============================================================

\paragraph{Quantum feature maps and kernel methods.}
The idea of using quantum circuits as feature maps for supervised learning
was formalised by \citet{havlicek2019}, who demonstrated that quantum-enhanced
feature spaces can express kernels that are classically hard to evaluate.
\citet{schuld2019} showed that variational quantum circuits are equivalent to
kernel methods operating in a quantum feature Hilbert space, providing a
theoretical bridge between quantum computing and classical kernel theory.
\citet{perez2020} extended this line of work with a data re-uploading strategy
that enables a single qubit to act as a universal classifier by encoding inputs
multiple times through the circuit.

\paragraph{Dequantisation and classical matching.}
A growing body of work has shown that classical algorithms can match or
approximate quantum methods under certain conditions.
\citet{tang2019} introduced the technique of dequantisation, demonstrating
that a classical algorithm can replicate the polynomial speedup of a quantum
recommendation system when given sampling access to the data.
\citet{huang2021} showed that, given sufficient data, classical machine
learning models can match the predictive performance of quantum kernel methods
on benchmark tasks, raising the question of whether quantum advantage in ML
requires a data regime beyond current experiments.
\citet{kubler2021} characterised the inductive bias of quantum kernels and
showed that it is governed by the geometry of the data distribution, implying
that classical kernels with matching geometry can replicate quantum kernel
behaviour.

\paragraph{Limitations and failure modes of QML.}
Several theoretical results constrain the practical reach of quantum ML models.
\citet{cerezo2021} survey variational quantum algorithms and document the
barren plateau phenomenon, in which gradients vanish exponentially with system
size, making training intractable at scale.
\citet{abbas2021} analyse the effective dimension and Fisher information of
quantum neural networks, finding that while quantum models can exhibit high
effective dimension, this does not straightforwardly translate to improved
generalisation.
The concentration of quantum kernel values, whereby Gram matrix entries become
exponentially similar with increasing qubit count, further limits the
discriminative power of quantum kernels on large
problems~\citep{thanasilp2022}.

\paragraph{Explicit feature maps in classical ML.}
The strategy of computing an explicit feature map to approximate a kernel
was introduced by \citet{rahimi2007} through Random Fourier Features (RFF),
which approximate the RBF kernel via random cosine projections and serve as
one of the strongest classical baselines in this benchmark.
The broader framework of kernel methods and reproducing kernel Hilbert spaces
that underlies both classical and quantum feature map approaches is surveyed
in \citet{scholkopf2002}.

\paragraph{Positioning of this work.}
Prior comparisons between QIE and classical methods have focused primarily
on accuracy endpoints.
This work differs by providing a mechanistic spectral attribution, in terms
of effective rank, condition number, and CKA, that explains \emph{why} each
encoding succeeds or fails rather than simply reporting whether it does.
The benchmark is also the first to evaluate all three canonical encodings
(amplitude, angle, basis) jointly across ten diverse datasets under a strictly
matched representational budget.
